\section{Go、Bitter Lesson 与通用 Inference Compute}

围棋提供了第二种对照。它是完全信息游戏，没有隐藏手牌，却仍依赖 search 将 policy/value network 转化为强行动。讲者借 AlphaGo Zero 的 ablation 和 Bitter Lesson，把“搜索”从特定博弈算法扩展为一个更宽的定义：在收到输入以后继续投入计算，以改善这一次输出。

\subsection{AlphaGo Zero：网络不是搜索的替代品}

图中的柱状条按系统版本展示棋力。Full AlphaGo Zero 同时使用训练好的网络和 Monte Carlo Tree Search；移除 MCTS 后，即使网络本身来自强 self-play，表现仍会大幅下降到 expert-human 以下。这里并不是说 MCTS 可直接搬到所有任务，而是说明训练得到的启发式与推理时展开彼此互补。

\lecturefigure{slide-06-search-go.jpg}{AlphaGo Zero 的搜索消融：Full system 与 no-search 版本差距明显}{官方录像恢复幻灯片 V006；Silver et al., Science 2017}

\begin{knowledgebox}{读图：柱高是棋力，不是训练 loss}
这张图比较的是最终对局水平。policy/value network 提供 move prior 和 leaf evaluation，MCTS 把它们组合成更可靠的局面决策。移除搜索后，模型并非“什么都不会”，但无法把局部预测组织成足够强的长程行动序列。因此不能把模型参数和 search visits 当成可完全互换的单一尺度。
\end{knowledgebox}

\subsection{课堂中的 100,000 倍类比应怎样理解}

讲者用一个非常粗的规则帮助建立数量级直觉：若 base system 的某项规模每扩大十倍，只换来约 50 Elo，而 search 带来约 250 Elo，那么达到同样 Elo 差距需要五次十倍扩张，即
\begin{equation}
R_{\text{equiv}}
\approx 10^{\Delta E/50}
=10^{250/50}=10^5.
\end{equation}
$\Delta E$ 是待替代的 Elo 增益，50 是课堂中的经验斜率，$R_{\text{equiv}}$ 是粗略规模倍数。这个计算的价值在于提醒我们不要忽视大幅曲线位移，而不是宣称 search 在任意任务上都固定等价于十万倍训练计算。

\teachervoice{老师明确把“约 100,000 倍”当作便于讨论的 extrapolation，而不是定律。真正主张是：如果只沿着 base-model scaling 轴优化，我们可能会错过一个便宜得多、但需要新系统设计的 inference axis。}

\subsection{Bitter Lesson：一般方法与可扩展计算}

Richard Sutton 的 Bitter Lesson 认为，长期最可靠的 AI 进步往往来自能利用不断增长计算的一般方法，尤其是 learning 与 search，而不是把大量人类知识硬编码进系统。课堂引用它，是为了把扑克和围棋放在同一条历史线上：领域知识仍有用，但应当服务于可扩展的计算过程，而不是替代它。

\lecturefigure{slide-07-bitter-lesson.jpg}{The Bitter Lesson：search 与 learning 是可随计算扩展的一般方法}{官方录像恢复幻灯片 V007；Richard Sutton, 2019}

\begin{warningbox}{“一般方法”不等于“无领域结构”}
Libratus、Pluribus、AlphaGo Zero 与 Cicero 都包含大量领域建模。Bitter Lesson 在本讲中的作用是选择优化方向：优先建设能吃到更多计算的机制，而不是声称所有 hand-crafted representation 都应立刻删除。
\end{warningbox}

\subsection{下一目标：能否做 general inference computation}

第八张 slide 把问题直接推向语言模型：如果 base model 已经很大，是否可以用比训练便宜得多的 inference compute 继续提升单次输出；如果某个答案价值极高，是否应该允许系统花几分钟、几小时甚至更久。这里的 \term{search} 是广义词，不限于显式树结构，也可以是反复生成、验证、修正或调用环境反馈。

\lecturefigure{slide-08-next-goal-generality.jpg}{从博弈搜索走向通用 inference-time computation}{官方录像恢复幻灯片 V008；Stanford CS25 Lecture 14}

\begin{importantbox}{价值敏感的推理预算}
合理的 inference budget 不应只由输入长度决定，还应由错误代价、可验证性和机会成本决定。推荐一首歌可以快速回答；药物设计、芯片布线或关键谈判则值得生成更多候选并进行更严格评估。系统目标是最大化“额外计算带来的边际价值”，而不是统一让所有请求变慢。
\end{importantbox}

\teachervoice{讲者把 chain-of-thought 视为一种很通用、但仍有限的额外计算方式：模型通过生成中间 token 展开问题，却不一定拥有可回溯、可验证或能长期维护状态的规划器。课堂的开放问题是：是否存在比 CoT 更一般、更可靠的 planning mechanism。}

\subsection{本章小结}

Go 的消融证明强 network 仍需要 search，Bitter Lesson 提供了历史解释，而 generality slide 则把研究议程转向可变预算的推理系统。接下来，Diplomacy 会把难度再提高一层：环境中的其他“状态”是会说话、会合作、也会改变主意的人。

